% ============================================================
% 3.2.3.2 Employment and Unemployment — Mark Scheme
% AQA AS Economics
% ============================================================

\documentclass[12pt,a4paper]{article}
\usepackage[margin=1in,headheight=15pt]{geometry}
\usepackage{enumitem}
\usepackage{amsmath}
\usepackage{fancyhdr}
\usepackage{booktabs}
\usepackage{array}
\RequirePackage{tcolorbox}
\tcbuselibrary{skins}

% Key Point — yellow; for the single most important idea on a slide
\newtcolorbox{keypoint}{
    enhanced,
    colback=yellow!12,
    colframe=yellow!60!black,
    fonttitle=\bfseries,
    title=Key Point,
    arc=2pt,
    boxrule=0.8pt,
}

% Formula — blue; for named equations and their variable key
\newtcolorbox{formula}[1][Formula]{
    enhanced,
    colback=blue!5,
    colframe=blue!45!black,
    fonttitle=\bfseries,
    title=#1,
    arc=2pt,
    boxrule=0.8pt,
}

% Example Question — teal; for exam-style questions posed to students
\newtcolorbox{examplequestion}[1][Example Question]{
    enhanced,
    colback=teal!6,
    colframe=teal!55!black,
    fonttitle=\bfseries,
    title=#1,
    arc=2pt,
    boxrule=0.8pt,
}

% Example Solution — green; always follows an examplequestion block
\newtcolorbox{examplesolution}[1][Solution]{
    enhanced,
    colback=green!6,
    colframe=green!45!black,
    fonttitle=\bfseries,
    title=#1,
    arc=2pt,
    boxrule=0.8pt,
}

% ---- Worksheet environments --------------------------------

% Scenario — blue; for group discussion prompts on worksheets
% Usage: \begin{scenario}{Scenario 1: Title with commas, fine} ... \end{scenario}
\newtcolorbox{scenario}[1]{
    enhanced,
    colback=blue!5,
    colframe=blue!45!black,
    fonttitle=\bfseries,
    title={#1},
    arc=2pt,
    boxrule=0.8pt,
}

% Instructions — blue; for worksheet-level instruction blocks
\newtcolorbox{instructions}{
    enhanced,
    colback=blue!5,
    colframe=blue!45!black,
    fonttitle=\bfseries,
    title=Instructions,
    arc=2pt,
    boxrule=0.8pt,
}

% Extract — grey; for data response source material
% Usage: \begin{extract}{Source: Author, Year} ... \end{extract}
\newtcolorbox{extract}[1]{
    enhanced,
    colback=gray!8,
    colframe=gray!50,
    fonttitle=\bfseries,
    title={#1},
    arc=2pt,
    boxrule=0.8pt,
}


\pagestyle{fancy}
\fancyhf{}
\lhead{Dr Jac Thomas \quad\textit{Mark Scheme}}
\rhead{Topic index: 3.2.3.2}
\rfoot{Page \thepage}

% Award marks generously for equivalent correct answers.
% Do not penalise candidates for poor spelling or grammar
% unless it obscures the economic meaning.

\begin{document}

\begin{center}
    \Large\textbf{Employment and Unemployment --- Mark Scheme}\\[5pt]
    \normalsize Topic index: 3.2.3.2 \quad|\quad AQA AS Economics\\[3pt]
    \textit{Total Marks: 60}
\end{center}

\vspace{1em}

% ============================================================
\section*{Section A: Key Terms (10 marks)}
% ============================================================

\textit{Award 1 mark for a basic correct statement; 2 marks for a precise, complete definition.}

\vspace{0.5em}

\begin{enumerate}[leftmargin=*, label=\textbf{\arabic*.}]

    \item \textbf{Unemployment} \hfill (2 marks)\\
    People of working age who are without a job but are actively seeking work and are available to start work. \\
    \textit{1 mark: without a job / not in employment. 2 marks: must include actively seeking work (or equivalent).}

    \vspace{0.5em}

    \item \textbf{The Claimant Count} \hfill (2 marks)\\
    A measure of unemployment that counts the number of people claiming unemployment-related benefits (Universal Credit or Jobseeker's Allowance).\\
    \textit{1 mark: counts benefit claimants. 2 marks: must specify Universal Credit / JSA or equivalent.}

    \vspace{0.5em}

    \item \textbf{The Labour Force Survey (LFS)} \hfill (2 marks)\\
    A quarterly household survey that measures unemployment using the International Labour Organisation (ILO) definition: those without work, available to start work within two weeks, and actively seeking work in the past four weeks.\\
    \textit{1 mark: a survey-based measure. 2 marks: must reference ILO/three-part definition or equivalent.}

    \vspace{0.5em}

    \item \textbf{Structural unemployment} \hfill (2 marks)\\
    Unemployment caused by a long-run decline in demand for a particular type of labour, typically due to changes in technology, international competition, or a fall in demand for an industry's output, resulting in a mismatch between the skills of workers and the skills demanded by employers.\\
    \textit{1 mark: mismatch of skills / industry decline. 2 marks: must explain the cause (technology / structural change / deindustrialisation).}

    \vspace{0.5em}

    \item \textbf{Cyclical unemployment} \hfill (2 marks)\\
    Unemployment caused by a deficiency of aggregate demand during a downturn or recession in the economic cycle; also called demand-deficient unemployment. Output falls below full-employment output, creating a negative output gap.\\
    \textit{1 mark: linked to recession / falling demand. 2 marks: must reference aggregate demand deficiency or negative output gap.}

\end{enumerate}

\newpage

% ============================================================
\section*{Section B: Group Discussion --- Indicative Answers}
% ============================================================

\textit{No marks awarded here. Use these notes to guide class feedback.}

\vspace{0.5em}

\textbf{Scenario 1: Tata Steel, Port Talbot}
\begin{itemize}
    \item Type: \textbf{Structural unemployment} --- the steel industry is in long-run decline; workers' skills are sector-specific.
    \item Cause: Deindustrialisation, cheaper global competition, and the shift away from heavy industry.
    \item Policy: Supply-side --- retraining and reskilling programmes (e.g.\ apprenticeships, vocational courses), relocation subsidies to improve geographical mobility, regional investment to attract new industries to Port Talbot.
\end{itemize}

\vspace{0.5em}

\textbf{Scenario 2: The 2008 Financial Crisis}
\begin{itemize}
    \item Type: \textbf{Cyclical unemployment} --- caused by a collapse in aggregate demand.
    \item Cause: Credit crunch reduced consumer spending and business investment; AD fell sharply.
    \item Policy: Demand-side --- expansionary fiscal policy (increased government spending, tax cuts) or monetary policy (lower interest rates, quantitative easing) to boost AD and close the negative output gap.
\end{itemize}

\vspace{0.5em}

\textbf{Scenario 3: UK Coastal Tourism}
\begin{itemize}
    \item Type: \textbf{Seasonal unemployment} --- predictable variation in labour demand across the year.
    \item Cause: Demand for tourism services is concentrated in summer months.
    \item Discussion: Not necessarily a major problem --- workers may plan for it or hold multiple seasonal jobs. Policy options are limited and may not be worth the cost; encouraging year-round tourism (e.g.\ winter events) could help but cannot fully eliminate the pattern.
\end{itemize}

\newpage

% ============================================================
\section*{Section C: Multiple Choice (10 marks, 1 each)}
% ============================================================

\vspace{0.5em}

\begin{table}[h]
    \centering
    \renewcommand{\arraystretch}{1.3}
    \begin{tabular}{@{}ccc@{}}
        \toprule
        \textbf{Question} & \textbf{Answer} & \textbf{Correct option} \\ \midrule
        1  & C & Claiming Universal Credit or Jobseeker's Allowance \\
        2  & B & Without work, available for work, and actively seeking work \\
        3  & C & Structural unemployment \\
        4  & C & Frictional unemployment \\
        5  & B & A rise in income tax reducing consumer spending and aggregate demand \\
        6  & B & Making UK exports cheaper and more competitive abroad \\
        7  & B & Spare productive capacity and above-average unemployment \\
        8  & C & Introducing retraining schemes to improve occupational mobility \\
        9  & A & Lower UK export demand, reducing output and employment \\
        10 & B & Not all unemployed people are eligible for or choose to claim benefits \\ \bottomrule
    \end{tabular}
\end{table}

\newpage

% ============================================================
\section*{Section D: Quantitative Skills (8 marks)}
% ============================================================

\begin{enumerate}[leftmargin=*, label=\textbf{\arabic*.}]

    \item \textbf{LFS unemployment rate for 2023} \hfill (2 marks)\\
    \[
        \text{Unemployment rate} = \frac{\text{Number unemployed}}{\text{Labour force}} \times 100
        = \frac{1.40}{35.00} \times 100 = \mathbf{4.0\%}
    \]
    \textit{1 mark: correct formula / method. 1 mark: correct answer (4.0\%).}

    \vspace{0.5em}

    \item \textbf{Number unemployed (LFS) in 2024} \hfill (1 mark)\\
    \[
        35.00 - 33.25 = \mathbf{1.75 \text{ million}}
    \]
    \textit{1 mark: 1.75 million (or 1,750,000).}

    \vspace{0.5em}

    \item \textbf{LFS unemployment rate for 2024} \hfill (2 marks)\\
    \[
        \frac{1.75}{35.00} \times 100 = \mathbf{5.0\%}
    \]
    \textit{1 mark: correct method using answer from Q2. 1 mark: correct answer (5.0\%). Award follow-through marks if Q2 answer is wrong but method is correct.}

    \vspace{0.5em}

    \item \textbf{Percentage change in unemployment, 2023--2024} \hfill (2 marks)\\
    \[
        \frac{1.75 - 1.40}{1.40} \times 100 = \frac{0.35}{1.40} \times 100 = \mathbf{25\%}
    \]
    \textit{1 mark: correct formula. 1 mark: correct answer (25\%). Follow-through from Q2 if applicable.}

    \vspace{0.5em}

    \item \textbf{Claimant count as \% of LFS, 2023} \hfill (1 mark)\\
    \[
        \frac{0.98}{1.40} \times 100 = \mathbf{70\%}
    \]
    \textit{1 mark: 70\%.}

\end{enumerate}

\newpage

% ============================================================
\section*{Section E: Diagrams (10 marks)}
% ============================================================

\textit{Mark for: correct shape of curves, correct labelling, correct equilibrium position, and correct identification of the output gap (Q1) or shift (Q2).}

\vspace{0.5em}

\textbf{Question 1 --- Cyclical unemployment} \hfill (5 marks)

\begin{itemize}
    \item \textbf{[1]} Downward-sloping AD curve, upward-sloping SRAS curve, vertical LRAS curve --- all labelled.
    \item \textbf{[1]} AD and SRAS intersect \textit{to the left of} LRAS (i.e.\ at output $Y_e < Y_f$).
    \item \textbf{[1]} $Y_e$ (actual output) and $Y_f$ (full-employment output) both labelled on x-axis.
    \item \textbf{[1]} Negative output gap clearly indicated between $Y_e$ and $Y_f$ (e.g.\ bracket, arrow, or label).
    \item \textbf{[1]} Equilibrium price level $P_e$ labelled on y-axis.
\end{itemize}

\vspace{0.8em}

\textbf{Question 2 --- AD increase closing the output gap} \hfill (5 marks)

\begin{itemize}
    \item \textbf{[1]} Initial position reproduced correctly (LRAS, SRAS, AD$_1$, $E_1$ labelled).
    \item \textbf{[1]} AD shifts rightward, labelled AD$_2$.
    \item \textbf{[1]} New equilibrium $E_2$ at or closer to $Y_f$.
    \item \textbf{[1]} New (higher) price level $P_2$ labelled.
    \item \textbf{[1]} Reduction in output gap shown or stated --- unemployment falls as output rises toward $Y_f$.
\end{itemize}

\newpage

% ============================================================
\section*{Section F: Short Response (7 marks)}
% ============================================================

\begin{enumerate}[leftmargin=*, label=\textbf{\arabic*.}]

    \item \textbf{Two differences between the claimant count and the LFS} \hfill (4 marks)\\
    \textit{Award 2 marks per difference: 1 mark for identifying the difference, 1 mark for developing it.}

    \vspace{0.3em}
    Indicative differences (any two):
    \begin{itemize}
        \item \textbf{Coverage:} The claimant count only includes those claiming Universal Credit / JSA, whereas the LFS counts all who meet the ILO three-part definition regardless of benefit eligibility. The claimant count therefore tends to be lower.
        \item \textbf{Measurement method:} The claimant count is an administrative count of benefit records; the LFS is a quarterly household survey of around 40,000 households. The LFS is therefore subject to sampling error; the claimant count is not.
        \item \textbf{Definition of unemployment:} The claimant count uses eligibility for benefits as its criterion; the LFS uses the ILO definition (without work, available, actively seeking). Some people who meet the ILO definition may not claim benefits (e.g.\ those whose partner earns too much), and some claimants may not meet the ILO definition.
        \item \textbf{Sensitivity to policy:} The claimant count can change when benefit rules change (e.g.\ the introduction of Universal Credit), not just when labour market conditions change. The LFS is less affected by policy changes.
    \end{itemize}

    \vspace{0.5em}

    \item \textbf{One supply-side cause of rising unemployment} \hfill (3 marks)\\
    \textit{1 mark: identify the factor. 2 marks: explain the mechanism clearly.}

    \vspace{0.3em}
    Indicative answers (any one):
    \begin{itemize}
        \item \textbf{Skills mismatch / lack of human capital:} If workers lack the skills demanded by employers (e.g.\ due to technological change or poor education), vacancies go unfilled even while unemployment rises. This is structural unemployment --- it cannot be solved by boosting demand.
        \item \textbf{Geographical immobility:} Workers may be unable or unwilling to move to areas where jobs are available (e.g.\ due to housing costs or family ties), leaving unemployment high in some regions despite job vacancies elsewhere.
        \item \textbf{Occupational immobility:} Workers with highly specialised skills in a declining industry cannot easily retrain for growing sectors, causing long-term structural unemployment.
        \item \textbf{High minimum wage / labour costs:} If the minimum wage is set above the market-clearing wage, the quantity of labour demanded may fall below the quantity supplied, causing unemployment (classical unemployment). \textit{Note: this is contested --- accept if well reasoned.}
    \end{itemize}

\end{enumerate}

\newpage

% ============================================================
\section*{Section G: Essay (15 marks)}
% ============================================================

\textit{``To what extent do the appropriate policies for reducing unemployment depend on the type of unemployment?''}

\vspace{0.5em}

\begin{extract}{Marking guidance}
    Level 4 (13--15): Thorough analysis of the link between unemployment types and policies. Strong evaluation --- considers both where policy \textit{must} be tailored to cause and where a general approach (e.g.\ boosting AD) helps across types. Well-structured, precise terminology.\\[0.4em]
    Level 3 (9--12): Good analysis of at least two types and their appropriate policies. Some evaluation of whether policies are type-specific or general. May lack depth in one area.\\[0.4em]
    Level 2 (5--8): Some relevant analysis of types and/or policies, but links between them may be underdeveloped. Limited evaluation.\\[0.4em]
    Level 1 (1--4): Basic or largely descriptive answer. May list types or policies without linking them. Little or no evaluation.
\end{extract}

\vspace{0.8em}

\textbf{Indicative content}

\vspace{0.3em}
\textbf{Where type determines policy:}
\begin{itemize}
    \item \textbf{Structural unemployment} requires supply-side policies --- retraining programmes, improving occupational and geographical mobility, investment in education. Boosting AD will not resolve a skills mismatch; it may only create inflationary pressure.
    \item \textbf{Frictional unemployment} can be reduced by improving job-matching information (e.g.\ better job centres, online platforms) and reducing search frictions. Demand-side policy has little effect here since vacancies exist.
    \item \textbf{Seasonal unemployment} is largely unavoidable and may not warrant intervention; policy might focus on encouraging year-round industries in affected regions.
\end{itemize}

\vspace{0.3em}
\textbf{Where a general approach works:}
\begin{itemize}
    \item \textbf{Cyclical unemployment} responds directly to expansionary fiscal or monetary policy --- cutting interest rates or increasing government spending raises AD, closes the negative output gap, and increases employment across all sectors simultaneously.
    \item Some supply-side policies (e.g.\ cutting employers' National Insurance contributions) reduce the cost of hiring across all unemployment types.
\end{itemize}

\vspace{0.3em}
\textbf{Evaluation:}
\begin{itemize}
    \item In practice, multiple types of unemployment coexist and it is difficult to distinguish between them, complicating policy design.
    \item Demand-side expansion can reduce cyclical unemployment but risks inflation if applied when structural unemployment is the main problem (economy may already be near $Y_f$ in the affected sectors).
    \item Supply-side policies have long time lags (e.g.\ retraining takes years); cyclical unemployment may worsen before structural policies take effect.
    \item Trade-off: policies to reduce one type may worsen another (e.g.\ higher benefits to support structural unemployment may increase frictional unemployment by reducing job-search urgency).
    \item \textbf{Conclusion:} The type of unemployment is a necessary but not sufficient guide to policy choice. Cyclical unemployment is relatively amenable to general demand management; structural and frictional unemployment require targeted supply-side intervention. Identifying the correct type before prescribing policy is therefore essential.
\end{itemize}

\vfill
\begin{center}
    \textbf{[Total: 60 marks]}
\end{center}

\end{document}
